\clearpage
\VerdierPhysicalPageAnchor{0202}%
\begingroup\small\setlength{\parskip}{0pt}\setlength{\emergencystretch}{1.5em}
\setlength{\abovedisplayskip}{4pt plus 2pt minus 2pt}%
\setlength{\belowdisplayskip}{4pt plus 2pt minus 2pt}%

\VerdierSemanticAnchor{ch3d-2-3-3-proof-q-r-step}
\[
\alpha\colon L^{-r+1}\longrightarrow Z^{-r+1},
\qquad\text{où }L^{-r+1}\text{ est un objet de }M.
\]
Posons alors
$X_{q,r}=X_{q,r-1}(\alpha,-r+1)$ et
$v_{q,r}=p(\alpha,-r+1)$ (2.1.1). Il est clair que $X_{q,r}$ et
$v_{q,r}$ possèdent les propriétés $I')_r$,
$III')_r$, $IV')_r$. De plus, par définition du
complexe $X_{q,r}$ (2.1.1), on a une suite exacte :
\[
0\longrightarrow (X_{q,r})^{-r}\longrightarrow
L^{-r}\oplus L^{-r+1}\longrightarrow Z^{-r+1}\longrightarrow0.
\]
Il résulte alors de la propriété $(E_4)$ que $(X_{q,r})^{-r}$ est un
élément de $M'$ et de (2.2.6) que :
\[
M\text{-}\dim((X_{q,r})^{-r})
=M\text{-}\dim(Z^{-r+1})-1
=M\text{-}\dim(Y^q)-r-q,
\]
d'où la propriété $II')_r$.

\VerdierSemanticAnchor{ch3d-2-3-3-proof-q-completion}
On peut donc poursuivre la construction. Posons alors
$X_{q,-q+M\text{-}\dim(Y^q)}=X_q$ et désignons par
$u_q\colon X_q\longrightarrow X_{q-1}$ le morphisme obtenu en composant
les $v_{q,l}$. Il est clair que $X_q$ et $u_q$ possèdent les propriétés
$1)_q$, $2)_q$, $3)_q$, $4)_q$, $5)_q$. Ceci achève la démonstration de
la première assertion de la proposition. Pour démontrer la deuxième
assertion, il suffit de remarquer que sous les hypothèses supplémentaires,
on peut prendre dans la construction précédente $X_{n_0}=Y$, pour un
$n_0$ convenable.

\VerdierSemanticAnchor{ch3d-2-3-4}
\medskip\noindent\textbf{Remarque 2.3.4.}\quad
Dans la construction du complexe $X_{n_0}$ (\emph{cf.} démonstration de la
proposition (2.3.3)), nous avons en fait démontré le résultat suivant :
soit $M\subset M'$ deux ensembles d'objets de $\mathcal A$ tels que
$0\in M$, $M'$ possède les propriétés $(E_1)$ et $(E_4)$ et que tout
objet de $M'$ soit quotient d'un objet de $M$. Soit
$Y\in\operatorname{Ob}\operatorname{comp}^{-}(\mathcal A)$ un complexe
dont tous les composants sont des éléments de $M'$. Il existe un complexe
$X\in\operatorname{Ob}\operatorname{comp}^{-}(\mathcal A)$ dont tous les
composants sont des éléments de $M$ et un quasi-isomorphisme
$u\colon X\longrightarrow Y$.

\VerdierSemanticAnchor{ch3d-2-3-5}
\medskip\noindent\textbf{2.3.5.}\quad
Nous laissons au lecteur le soin d'énoncer les résultats qu'on obtient en
appliquant les résultats de ce numéro à la catégorie opposée à $\mathcal A$.
On dit qu'un ensemble $M$ d'objets de $\mathcal A$ possède la propriété
$(E_i)^o$, $1\leq i\leq6$, si l'ensemble $M$ d'objets de $\mathcal A^o$
possède la propriété $(E_i)$ de (2.2.1).

\VerdierSemanticAnchor{ch3d-2-4}
\setcounter{subsection}{3}%
\subsection{Applications.}

\VerdierSemanticAnchor{ch3d-2-4-1}
\noindent\textbf{Proposition 2.4.1.}\quad
\textit{a) Soient $\mathcal A$ une catégorie abélienne et $M$ un ensemble
d'objets de $\mathcal A$ possédant la propriété $(E_1)$ (2.2.1). La
catégorie $K^{-}(M)$ (1.1.4) est
$\operatorname{Ac}(\mathcal A)$-localisante à gauche (chap. II, 2.3.5) si
et seulement si l'ensemble $M$ possède les propriétés $(E_2)$ et $(E_3)$.}

\noindent\textit{b) Soient $M$ un ensemble d'objets possédant les
propriétés $(E_1)$, $(E_2)$ et $(E_3)$ et $M^{\wedge}$ l'ensemble associé
à $M$, défini en (2.2.1). Le foncteur naturel :}
\[
K^{-}(M)/\operatorname{Ac}^{-}(M)
\longrightarrow
K^{-}(M^{\wedge})/\operatorname{Ac}^{-}(M^{\wedge})
\]
\endgroup
